\chapter{Analyse et spécification des besoins}
\label{chap:analyse}

Les deux chapitres précédents ont établi ce qu'il fallait construire et ce que la
science savait déjà faire. Celui-ci formalise la commande. Nous identifions d'abord
les acteurs du système et leurs responsabilités, puis nous exprimons les besoins
fonctionnels et non fonctionnels. Ces besoins sont ensuite traduits en un
\emph{Product Backlog} de soixante récits utilisateur, priorisés et estimés, dont
l'ordonnancement produit le découpage du projet en six sprints. Le chapitre se clôt
sur la modélisation des cas d'utilisation.

% =============================================================================
\section{Identification des acteurs}
% =============================================================================

Un acteur désigne toute entité qui interagit avec le système sans en faire partie. On
distingue les acteurs \emph{principaux}, qui déclenchent les cas d'utilisation pour en
retirer une valeur, des acteurs \emph{secondaires}, sollicités par le système pour
mener une opération à bien. Le tableau~\ref{tab:acteurs} recense les quatre acteurs
de \hypnoria{}.

\begin{table}[H]
\centering
\caption{Acteurs du système et leurs responsabilités}
\label{tab:acteurs}
\small
\begin{tabular}{@{}L{3.0cm}cL{8.2cm}@{}}
\toprule
\textbf{Acteur} & \textbf{Type} & \textbf{Responsabilités} \\
\midrule
Médecin du sommeil & principal &
Constituer et suivre son registre de patients, lancer l'analyse d'un enregistrement,
interpréter l'hypnogramme et la prédiction de sévérité, annoter, produire le compte
rendu, solliciter un confrère \\
\addlinespace[3pt]
Administrateur & principal &
Créer et administrer les comptes médecins, gérer les échéances de licence, rattacher
les praticiens à un établissement, superviser l'activité et consulter le journal
d'audit \\
\addlinespace[3pt]
Moteur d'inférence & secondaire &
Charger les modèles, prétraiter le signal, produire la séquence de stades, calculer
les métriques et les attributions d'explication \\
\addlinespace[3pt]
Service de stockage & secondaire &
Recevoir les fichiers volumineux, les conserver de façon durable et délivrer des liens
d'accès temporaires \\
\bottomrule
\end{tabular}
\end{table}

Deux remarques sur ce découpage. Le \textbf{moteur d'inférence} est modélisé comme un
acteur distinct bien qu'il réside dans le même processus que l'API : ce choix
matérialise le fait qu'il pourrait être déporté sur une machine dédiée sans modifier
les cas d'utilisation. Le \textbf{service de stockage}, lui, est réellement externe.

Le rôle d'administrateur d'établissement, prévu dans le modèle de données, n'a pas été
implémenté : il figure parmi les perspectives.

% =============================================================================
\section{Besoins fonctionnels}
% =============================================================================

Les besoins fonctionnels ont été organisés en neuf \emph{épics}, regroupements
thématiques de fonctionnalités partageant un même objectif métier. Cette structure,
présentée au tableau~\ref{tab:epics}, sert de colonne vertébrale au backlog.

\begin{table}[H]
\centering
\caption{Regroupement des besoins fonctionnels en épics}
\label{tab:epics}
\small
\begin{tabular}{@{}clL{7.4cm}cc@{}}
\toprule
\textbf{Épic} & \textbf{Intitulé} & \textbf{Portée} & \textbf{Récits} &
\textbf{Points} \\
\midrule
E1 & Authentification & Connexion, jetons, hachage, contrôle des rôles, application des
licences & 5 & 18 \\
\addlinespace[2pt]
E2 & Administration & Comptes médecins, cycle de vie, établissements, audit,
statistiques & 10 & 45 \\
\addlinespace[2pt]
E3 & Dossier patient & Registre des patients, examens, historique, tableau de bord
& 6 & 24 \\
\addlinespace[2pt]
E4 & Classification & Lecture EDF, prétraitement, modèles de base, métriques AASM
& 12 & 95 \\
\addlinespace[2pt]
E5 & Sévérité du SAOS & Empilement, extraction de marqueurs, prédiction, explications
& 8 & 65 \\
\addlinespace[2pt]
E6 & Visualisation & Hypnogramme, oxymétrie, annotation, normes de référence & 5 & 28 \\
\addlinespace[2pt]
E7 & Collaboration & Discussions rattachées à un examen, boîte de réception & 3 & 18 \\
\addlinespace[2pt]
E8 & Archivage et export & Stockage cloud, liens signés, compte rendu, exports
& 6 & 36 \\
\addlinespace[2pt]
E9 & Industrialisation & Conteneurisation, intégration continue, multilinguisme
& 5 & 29 \\
\midrule
\multicolumn{3}{@{}r}{\textbf{Total}} & \textbf{60} & \textbf{358} \\
\bottomrule
\end{tabular}
\end{table}

% =============================================================================
\section{Besoins non fonctionnels}
% =============================================================================

Les exigences non fonctionnelles conditionnent l'acceptabilité du système autant que
ses fonctions. Elles sont d'autant plus contraignantes ici que le domaine — la donnée
de santé — impose des obligations propres. Le tableau~\ref{tab:bnf} les énonce avec,
pour chacune, le moyen de vérification retenu : un besoin non fonctionnel qu'on ne
sait pas vérifier n'est qu'une intention.

\begin{table}[H]
\centering
\caption{Besoins non fonctionnels et moyens de vérification}
\label{tab:bnf}
\small
\begin{tabular}{@{}L{2.7cm}L{5.6cm}L{5.6cm}@{}}
\toprule
\textbf{Critère} & \textbf{Exigence} & \textbf{Vérification} \\
\midrule
Sécurité &
Aucun accès sans jeton valide ; mots de passe jamais stockés en clair ; cloisonnement
strict entre rôles &
Revue des dépendances de sécurité sur chaque route ; tentative d'accès croisé entre
comptes \\
\addlinespace[3pt]
Confidentialité &
Aucun fichier patient accessible par une URL publique ou devinable &
Vérification que tout lien délivré est signé et expire \\
\addlinespace[3pt]
Traçabilité &
Toute action sensible horodatée avec son auteur, journal exportable &
Inspection du journal après un scénario complet \\
\addlinespace[3pt]
Performance &
Analyse d'une nuit complète en quelques minutes ; interface non bloquée pendant les
transferts &
Chronométrage sur enregistrements réels ; transfert en arrière-plan observable \\
\addlinespace[3pt]
Interopérabilité &
Accepter les fichiers EDF quel que soit le constructeur ; importer et exporter en
formats ouverts &
Résolution des canaux sur des enregistrements de provenances différentes \\
\addlinespace[3pt]
Ergonomie &
Parcours d'analyse guidé, sans formation préalable ; retour visuel sur les traitements
longs &
Passage d'un utilisateur non initié sur le parcours complet \\
\addlinespace[3pt]
Multilinguisme &
Interface intégralement disponible en français et en anglais &
Revue de l'absence de chaînes non traduites \\
\addlinespace[3pt]
Portabilité &
Déploiement complet sur une machine vierge sans configuration manuelle &
Démarrage de la pile sur un poste neuf en une commande \\
\addlinespace[3pt]
Maintenabilité &
Séparation nette des responsabilités ; ajout d'un modèle sans toucher au reste &
Ajout d'une variante de modèle à titre d'essai \\
\bottomrule
\end{tabular}
\end{table}

% =============================================================================
\section{Product Backlog}
% =============================================================================

Le \emph{Product Backlog} est la liste ordonnée de tout ce que le produit doit offrir.
Chaque élément prend la forme d'un récit utilisateur suivant la formulation
canonique : \emph{en tant que} \dots, \emph{je veux} \dots, \emph{afin de} \dots. Cette
formulation force à énoncer le bénéficiaire et la finalité, et non seulement la
fonction.

Deux attributs complètent chaque récit. La \textbf{priorité} suit la méthode MoSCoW —
\emph{Must have} pour ce sans quoi le produit n'a pas de sens, \emph{Should have} pour
ce qui lui donne sa valeur, \emph{Could have} pour ce qui l'améliore sans le
conditionner. L'\textbf{estimation} s'exprime en points de complexité sur la suite de
Fibonacci, échelle relative qui compare les récits entre eux plutôt que de prétendre
prédire une durée.

Le backlog complet figure au tableau~\ref{tab:backlog}.

\newcommand{\epictitre}[1]{%
  \midrule
  \multicolumn{5}{@{}l}{\cellcolor{hypnolight}\textbf{#1}}\\
  \midrule}

\begin{longtable}{@{}p{1.25cm}L{9.1cm}cp{0.75cm}p{0.9cm}@{}}
\caption{Product Backlog --- 60 récits utilisateur, 358 points}
\label{tab:backlog}\\
\toprule
\textbf{ID} & \textbf{Récit utilisateur} & \textbf{Prio.} & \textbf{Pts} &
\textbf{Spr.} \\
\midrule
\endfirsthead

\multicolumn{5}{@{}l}{\small\itshape Tableau~\ref{tab:backlog} --- suite}\\
\toprule
\textbf{ID} & \textbf{Récit utilisateur} & \textbf{Prio.} & \textbf{Pts} &
\textbf{Spr.} \\
\midrule
\endhead

\midrule
\multicolumn{5}{@{}r}{\small\itshape suite page suivante}\\
\endfoot

\bottomrule
\endlastfoot

\small
US-01 & Me connecter avec mon adresse et mon mot de passe pour accéder à la plateforme
& M & 5 & S4 \\
US-02 & Consulter mon profil pour vérifier mon rôle et mon rattachement & M & 2 & S4 \\
US-03 & Stocker les mots de passe hachés pour qu'aucun ne soit lisible en base
& M & 3 & S4 \\
US-04 & Vérifier le rôle sur chaque route pour cloisonner médecins et administrateurs
& M & 5 & S4 \\
US-05 & Refuser la connexion d'un compte suspendu ou dont la licence a expiré
& M & 3 & S5 \\

\epictitre{E2 --- Administration de la plateforme}
US-06 & Créer un compte médecin pour ouvrir l'accès à un nouveau praticien & M & 5 & S5 \\
US-07 & Lister les médecins avec leur statut et leur échéance de licence & M & 3 & S5 \\
US-08 & Suspendre ou réactiver un compte et fixer sa date d'expiration & M & 5 & S5 \\
US-09 & Réinitialiser le mot de passe d'un médecin pour débloquer un accès perdu
& S & 3 & S5 \\
US-10 & Créer des établissements et y rattacher des médecins & S & 5 & S5 \\
US-11 & Journaliser chaque action sensible avec son auteur et son horodatage
& M & 5 & S5 \\
US-12 & Consulter le journal d'audit pour enquêter sur un incident & M & 5 & S5 \\
US-13 & Exporter le journal d'audit pour un contrôle de conformité & S & 3 & S5 \\
US-14 & Consulter les statistiques globales pour suivre l'adoption & S & 3 & S5 \\
US-15 & Visualiser la chaîne de traitement des modèles & C & 8 & S6 \\

\epictitre{E3 --- Dossier patient et examens}
US-16 & Enregistrer un patient avec son âge, son IMC et son sexe & M & 3 & S5 \\
US-17 & Lister mes patients avec leur historique d'examens & M & 5 & S5 \\
US-18 & Consulter la fiche détaillée d'un patient et tous ses examens & M & 3 & S5 \\
US-19 & Créer un examen rattaché à un patient pour y associer les résultats
& M & 5 & S5 \\
US-20 & Mettre à jour la sévérité et les résultats d'un examen & M & 3 & S5 \\
US-21 & Voir la répartition des risques d'apnée de ma patientèle & S & 5 & S6 \\

\epictitre{E4 --- Classification automatique des stades}
US-22 & Lire un fichier EDF et en extraire les dérivations physiologiques & M & 8 & S1 \\
US-23 & Résoudre les canaux malgré l'hétérogénéité des conventions de nommage
& M & 5 & S1 \\
US-24 & Rééchantillonner, segmenter en époques de 30 s et normaliser & M & 8 & S1 \\
US-25 & Valider la chaîne sur un jeu public en attendant l'accès aux cohortes
& M & 8 & S1 \\
US-26 & Prétraiter les enregistrements de la cohorte clinique à l'échelle & M & 8 & S2 \\
US-27 & Classer chaque époque avec un réseau convolutif multi-échelle & M & 13 & S2 \\
US-28 & Classer chaque époque avec un réseau récurrent à pooling attentionnel
& M & 13 & S2 \\
US-29 & Classer chaque époque avec un Transformer par patchs & M & 13 & S2 \\
US-30 & Choisir la configuration d'analyse selon le montage disponible & S & 5 & S4 \\
US-31 & Inspecter les canaux d'un fichier avant de lancer l'analyse & S & 3 & S4 \\
US-32 & Calculer les métriques AASM de l'enregistrement & M & 8 & S4 \\
US-33 & Conserver les modèles en mémoire entre deux requêtes & S & 3 & S4 \\

\epictitre{E5 --- Ensemble et prédiction de la sévérité}
US-34 & Combiner les trois modèles par un méta-apprenant & M & 13 & S3 \\
US-35 & Extraire les marqueurs d'architecture du sommeil depuis l'hypnogramme
& M & 8 & S3 \\
US-36 & Prédire la sévérité du SAOS en quatre classes & M & 13 & S3 \\
US-37 & Exposer les facteurs aggravants et protecteurs de la prédiction & M & 8 & S3 \\
US-38 & Saisir les données oxymétriques et les index d'éveil & M & 5 & S4 \\
US-39 & Lancer une prédiction sans fichier, à partir d'une saisie manuelle & S & 8 & S4 \\
US-40 & Importer un rapport externe pour préremplir le formulaire clinique
& S & 5 & S4 \\
US-41 & Lire une interprétation clinique automatique des anomalies détectées
& S & 5 & S4 \\

\epictitre{E6 --- Visualisation, annotation et normes}
US-42 & Visualiser l'hypnogramme de façon interactive & M & 8 & S4 \\
US-43 & Voir la saturation et la densité d'événements sur le même axe temporel
& S & 5 & S4 \\
US-44 & Suivre la progression de l'analyse en cours & C & 5 & S4 \\
US-45 & Annoter une époque précise pour consigner une observation & S & 5 & S6 \\
US-46 & Situer chaque métrique par rapport aux normes de référence & S & 5 & S6 \\

\epictitre{E7 --- Collaboration entre praticiens}
US-47 & Ouvrir une discussion sur un examen avec un confrère & S & 8 & S6 \\
US-48 & Échanger des messages dans un fil rattaché à l'examen & S & 5 & S6 \\
US-49 & Consulter une boîte de réception centralisée & S & 5 & S6 \\

\epictitre{E8 --- Archivage et exports}
US-50 & Archiver le fichier brut sans bloquer l'interface & M & 8 & S5 \\
US-51 & Archiver l'hypnogramme et le compte rendu clinique & S & 5 & S5 \\
US-52 & Délivrer des liens d'accès à durée limitée & M & 5 & S5 \\
US-53 & Exporter les marqueurs calculés pour un autre outil & S & 5 & S4 \\
US-54 & Composer le compte rendu en choisissant les sections à inclure & M & 8 & S6 \\
US-55 & Synchroniser l'hypnogramme annoté vers le dossier & S & 5 & S6 \\

\epictitre{E9 --- Industrialisation et confort d'usage}
US-56 & Conteneuriser les trois services pour un déploiement reproductible
& M & 8 & S6 \\
US-57 & Orchestrer la pile en une seule commande & M & 5 & S6 \\
US-58 & Automatiser tests et construction d'images à chaque livraison & S & 8 & S6 \\
US-59 & Basculer l'interface entre français et anglais & S & 5 & S6 \\
US-60 & Basculer en thème sombre pour la lecture nocturne & C & 3 & S6 \\
\end{longtable}

\noindent
La formulation \emph{en tant que \dots, je veux \dots, afin de \dots} a été abrégée
dans le tableau pour tenir sur une ligne ; le bénéficiaire se déduit de l'épic.

% =============================================================================
\section{Planification et déroulement des sprints}
% =============================================================================

\subsection{Le choix d'itérations de quatre semaines}

Les vingt-six semaines du stage ont été découpées en six sprints de quatre semaines,
suivis de deux semaines consacrées à la rédaction et à la préparation de la
soutenance. Ce rythme s'écarte des deux semaines habituelles pour une raison
technique : l'entraînement d'un modèle profond sur près de deux cent mille exemples
se compte en heures de calcul, et douze variantes devaient être produites. Une
itération de deux semaines n'aurait pas permis de livrer un incrément évaluable
pendant la phase de modélisation.

\subsection{Contenu des sprints}

Le tableau~\ref{tab:sprints} présente l'objectif, le contenu et le livrable de chaque
sprint.

\begin{table}[H]
\centering
\caption{Découpage du projet en six sprints}
\label{tab:sprints}
\small
\begin{tabular}{@{}L{1.5cm}L{2.5cm}L{5.0cm}L{4.4cm}c@{}}
\toprule
\textbf{Sprint} & \textbf{Période} & \textbf{Objectif} & \textbf{Livrable} &
\textbf{Pts} \\
\midrule
S1 & 26 jan.\newline 22 fév. &
Cadrer le projet, lancer la demande d'accès aux cohortes et prouver la faisabilité de
la chaîne sans attendre l'autorisation &
Chaîne de prétraitement opérationnelle et premiers modèles de faisabilité & 29 \\
\addlinespace[3pt]
S2 & 23 fév.\newline 22 mars &
Passer du jeu d'amorçage à la cohorte clinique et couvrir les trois familles
d'architectures &
Douze modèles de base entraînés et évalués & 47 \\
\addlinespace[3pt]
S3 & 23 mars\newline 19 avr. &
Faire coopérer les architectures, puis franchir le pas du staging vers le diagnostic
expliqué &
Quatre ensembles et le modèle de sévérité avec ses explications & 42 \\
\addlinespace[3pt]
S4 & 20 avr.\newline 17 mai &
Rendre les modèles consommables et donner au praticien un parcours guidé &
API documentée et parcours complet du dépôt au compte rendu & 80 \\
\addlinespace[3pt]
S5 & 18 mai\newline 14 juin &
Transformer l'outil d'analyse en plateforme gouvernée et traçable &
Dossier patient persistant, archivage cloud, espace d'administration & 77 \\
\addlinespace[3pt]
S6 & 15 juin\newline 12 juil. &
Couvrir le travail à plusieurs, la restitution documentaire et la reproductibilité &
Collaboration, compte rendu exportable, pile conteneurisée & 83 \\
\bottomrule
\end{tabular}
\end{table}

\subsection{Vélocité observée}

La vélocité mesure le nombre de points effectivement livrés par sprint. Sa
progression, représentée à la figure~\ref{fig:velocite}, est instructive.

\begin{figure}[H]
\centering
\begin{tikzpicture}[x=1cm,y=1cm]

  % Axe des ordonnées
  \draw[hypnogray!60,line width=0.6pt] (0,0) -- (0,4.6);
  \foreach \v/\y in {0/0, 20/1.0, 40/2.0, 60/3.0, 80/4.0} {
    \draw[hypnogray!60] (-0.1,\y) -- (0,\y);
    \node[left,font=\scriptsize,text=hypnogray] at (-0.12,\y) {\v};
    \draw[hypnogray!18,line width=0.4pt] (0,\y) -- (10.6,\y);
  }
  \node[rotate=90,font=\scriptsize,text=hypnogray] at (-1.0,2.3) {points livrés};

  % Barres
  \foreach \i/\lab/\pts in {0/S1/29, 1/S2/47, 2/S3/42, 3/S4/80, 4/S5/77, 5/S6/83} {
    \pgfmathsetmacro{\xa}{0.55 + \i*1.65}
    \pgfmathsetmacro{\xb}{\xa + 1.1}
    \pgfmathsetmacro{\hh}{\pts/20}
    \fill[hypnoblue!65] (\xa,0) rectangle (\xb,\hh);
    \draw[hypnoblue] (\xa,0) rectangle (\xb,\hh);
    \node[font=\scriptsize\bfseries,text=hypnoblue,anchor=south]
         at ({(\xa+\xb)/2},\hh+0.04) {\pts};
    \node[font=\scriptsize,text=hypnogray,anchor=north]
         at ({(\xa+\xb)/2},-0.06) {\lab};
  }

  % Axe des abscisses
  \draw[hypnogray!60,line width=0.6pt] (0,0) -- (10.6,0);

  % Vélocité moyenne
  \draw[hypnored,dashed,line width=0.8pt] (0,2.983) -- (10.6,2.983);
  \node[font=\scriptsize\itshape,text=hypnored,anchor=west] at (9.0,3.28)
       {moyenne : 60};

\end{tikzpicture}
\vspace{0.3cm}
\caption[Vélocité par sprint]%
        {Points de complexité livrés par sprint.}
\label{fig:velocite}
\end{figure}

Cette courbe ne décrit pas un plateau régulier, et c'est précisément ce qui la rend
crédible. Les deux premiers sprints sont absorbés par la montée en compétence sur le
cadre d'apprentissage profond et par des temps d'entraînement incompressibles : la
valeur livrée y est faible en points, mais elle conditionne tout le reste. Le sprint
S3 marque un creux relatif, la construction des ensembles supposant que les douze
modèles de base soient tous disponibles. À partir de S4, la chaîne est outillée et la
vélocité double, pour se maintenir jusqu'à la fin.

\subsection{Diagramme de Gantt}

La figure~\ref{fig:gantt} confronte la planification aux travaux réellement menés.

\begin{figure}[H]
\centering
\begin{ganttchart}[
    x unit                  = 0.375cm,
    y unit title            = 0.52cm,
    y unit chart            = 0.44cm,
    hgrid                   = {draw=hypnogray!18},
    vgrid                   = {*1{draw=hypnogray!18}},
    title/.append style     = {draw=hypnoblue!45, fill=hypnolight},
    title label font        = \tiny,
    title height            = 1,
    bar/.append style       = {draw=hypnoblue!85, fill=hypnoblue!55,
                               rounded corners=1pt},
    bar height              = 0.58,
    bar label font          = \scriptsize,
    bar label node/.append style = {text=hypnoblue!95},
    milestone/.append style = {fill=hypnored, draw=hypnored!85},
    milestone label font    = \scriptsize\itshape,
    milestone label node/.append style = {text=hypnored},
    milestone height        = 0.42,
  ]{1}{26}

  \gantttitle{Sprint 1}{4}
  \gantttitle{Sprint 2}{4}
  \gantttitle{Sprint 3}{4}
  \gantttitle{Sprint 4}{4}
  \gantttitle{Sprint 5}{4}
  \gantttitle{Sprint 6}{4}
  \gantttitle{Clôture}{2} \\
  \gantttitlelist{1,...,26}{1} \\

  \ganttbar{Étude bibliographique}{1}{3} \\
  \ganttbar{Demande d'accès NSRR}{1}{5} \\
  \ganttbar{Chaîne de prétraitement}{2}{4} \\
  \ganttbar{Amorçage sur jeu public}{3}{4} \\
  \ganttmilestone{Accès aux cohortes obtenu}{5} \\
  \ganttbar{Préparation de la cohorte}{5}{6} \\
  \ganttbar{Modèles de base}{6}{9} \\
  \ganttmilestone{12 modèles entraînés}{9} \\
  \ganttbar{Ensembles par empilement}{9}{10} \\
  \ganttbar{Modèle de sévérité}{10}{12} \\
  \ganttbar{Service d'inférence}{13}{15} \\
  \ganttbar{Interface d'analyse}{14}{17} \\
  \ganttbar{Base de données, archivage}{17}{19} \\
  \ganttbar{Administration, audit}{19}{20} \\
  \ganttbar{Collaboration, annotation}{21}{22} \\
  \ganttbar{Compte rendu, multilinguisme}{22}{23} \\
  \ganttmilestone{Plateforme complète}{22} \\
  \ganttbar{Conteneurisation, intégration}{23}{24} \\
  \ganttbar{Rédaction du mémoire}{21}{26} \\
  \ganttmilestone{Soutenance}{26}

\end{ganttchart}
\vspace{0.3cm}
\caption[Diagramme de Gantt du projet]%
        {Déroulement du projet sur les vingt-six semaines du stage.}
\label{fig:gantt}
\end{figure}

Le jalon le plus structurant est l'obtention de l'accès aux cohortes en semaine 5. Le
recouvrement volontaire entre la demande d'accès et les travaux d'amorçage sur le jeu
public traduit la décision évoquée au chapitre~\ref{chap:cadre} : ne pas subir un
délai administratif, mais l'occuper.

% =============================================================================
\section{Modélisation des cas d'utilisation}
% =============================================================================

\subsection{Diagramme de cas d'utilisation global}

La figure~\ref{fig:uc-global} présente la vue d'ensemble des interactions.

\begin{figure}[H]
\centering
\begin{tikzpicture}[
    x=1cm,y=1cm,
    uc/.style      = {ellipse, draw=hypnoblue!60, fill=hypnolight,
                      align=center, font=\scriptsize, text width=2.45cm,
                      inner sep=1.5pt, minimum height=0.92cm},
    lien/.style    = {hypnogray!75, line width=0.5pt},
    nomact/.style  = {font=\scriptsize\bfseries, align=center, text width=2.2cm},
  ]

  % ── Pictogramme d'acteur ────────────────────────────────────────────────
  \tikzset{acteur/.pic={
      \draw[hypnoblue,line width=0.8pt] (0,0) circle (0.17);
      \draw[hypnoblue,line width=0.8pt] (0,-0.17) -- (0,-0.66);
      \draw[hypnoblue,line width=0.8pt] (-0.28,-0.34) -- (0.28,-0.34);
      \draw[hypnoblue,line width=0.8pt] (0,-0.66) -- (-0.24,-1.04);
      \draw[hypnoblue,line width=0.8pt] (0,-0.66) -- (0.24,-1.04);}}

  % ── Frontière du système ────────────────────────────────────────────────
  \draw[hypnoblue!55,line width=0.9pt,rounded corners=3pt]
       (3.45,-6.40) rectangle (11.05,2.55);
  \node[font=\small\bfseries,text=hypnoblue,anchor=north] at (7.25,2.42)
       {Système \hypnoria{}};

  % ── Cas d'utilisation ───────────────────────────────────────────────────
  \node[uc] (auth)  at (7.25, 1.20) {S'authentifier};

  \node[uc] (pat)   at (5.35,-0.40) {Gérer ses patients};
  \node[uc] (ana)   at (5.35,-1.70) {Analyser un\\enregistrement};
  \node[uc] (osa)   at (5.35,-3.00) {Prédire la\\sévérité du SAOS};
  \node[uc] (annot) at (5.35,-4.30) {Annoter et\\exporter};
  \node[uc] (colab) at (5.35,-5.60) {Solliciter un\\confrère};

  \node[uc] (comp)  at (9.15,-0.40) {Gérer les comptes\\et licences};
  \node[uc] (etab)  at (9.15,-1.70) {Gérer les\\établissements};
  \node[uc] (audit) at (9.15,-3.00) {Consulter le\\journal d'audit};
  \node[uc] (stats) at (9.15,-4.30) {Suivre l'activité\\de la plateforme};

  % ── Acteurs principaux : libellé au-dessus, liaisons vers la droite ─────
  \pic at (1.15,0.00) {acteur};
  \node[nomact,anchor=south] at (1.15,0.42) {Médecin\\du sommeil};

  \pic at (13.35,0.00) {acteur};
  \node[nomact,anchor=south] at (13.35,0.42) {Administrateur};

  \foreach \c in {auth,pat,ana,osa,annot,colab}
    \draw[lien] (1.50,-0.20) -- (\c.west);

  \foreach \c in {auth,comp,etab,audit,stats}
    \draw[lien] (13.00,-0.20) -- (\c.east);

  % ── Acteurs secondaires : libellé en dessous ────────────────────────────
  \pic at (7.25,-7.55) {acteur};
  \node[nomact,anchor=north] at (7.25,-8.72) {Moteur\\d'inférence};

  \pic at (10.35,-7.55) {acteur};
  \node[nomact,anchor=north] at (10.35,-8.72) {Service de\\stockage};

  \draw[lien] (7.25,-7.28) -- (ana.east);
  \draw[lien] (7.25,-7.28) -- (osa.east);
  \draw[lien] (10.35,-7.28) -- (annot.east);

\end{tikzpicture}
\vspace{0.2cm}
\caption[Diagramme de cas d'utilisation global]%
        {Vue d'ensemble des cas d'utilisation de \hypnoria{}.}
\label{fig:uc-global}
\end{figure}

\subsection{Raffinement : analyse et prédiction}

Le cas « Analyser un enregistrement » est le c\oe ur fonctionnel du système. La
figure~\ref{fig:uc-analyse} le décompose et fait apparaître les relations
d'inclusion — sous-cas systématiquement exécutés — et d'extension — comportements
optionnels.

\begin{figure}[H]
\centering
\begin{tikzpicture}[
    x=1cm,y=1cm,
    uc/.style      = {ellipse, draw=hypnoblue!60, fill=hypnolight,
                      align=center, font=\scriptsize, text width=2.5cm,
                      inner sep=1.5pt, minimum height=0.92cm},
    ucp/.style     = {ellipse, draw=hypnored!75, fill=hypnored!8,
                      align=center, font=\scriptsize\bfseries, text width=2.5cm,
                      inner sep=1.5pt, minimum height=0.92cm},
    lien/.style    = {hypnogray!75, line width=0.5pt},
    incl/.style    = {-{Latex[length=2mm]}, hypnoblue!75, line width=0.5pt, dashed},
    ext/.style     = {-{Latex[length=2mm]}, hypnored!70, line width=0.5pt, dashed},
    et/.style      = {font=\tiny\itshape, fill=white, inner sep=1.2pt},
    nomact/.style  = {font=\scriptsize\bfseries, align=center, text width=1.8cm},
  ]

  \tikzset{acteur/.pic={
      \draw[hypnoblue,line width=0.8pt] (0,0) circle (0.17);
      \draw[hypnoblue,line width=0.8pt] (0,-0.17) -- (0,-0.66);
      \draw[hypnoblue,line width=0.8pt] (-0.28,-0.34) -- (0.28,-0.34);
      \draw[hypnoblue,line width=0.8pt] (0,-0.66) -- (-0.24,-1.04);
      \draw[hypnoblue,line width=0.8pt] (0,-0.66) -- (0.24,-1.04);}}

  % ── Frontière ───────────────────────────────────────────────────────────
  \draw[hypnoblue!55,line width=0.9pt,rounded corners=3pt]
       (2.75,-4.55) rectangle (12.60,3.15);
  \node[font=\small\bfseries,text=hypnoblue,anchor=north] at (7.65,3.02)
       {Module d'analyse};

  % ── Cas central et cas déclenchés par le médecin ────────────────────────
  \node[ucp] (ana) at (4.75, 0.75) {Analyser un\\enregistrement};
  \node[uc]  (exp) at (4.75,-2.20) {Exporter les\\marqueurs};
  \node[uc]  (pdf) at (4.75,-3.65) {Produire le\\compte rendu};

  % ── Sous-cas inclus ─────────────────────────────────────────────────────
  \node[uc] (chn) at (10.30, 2.10) {Choisir le montage\\et la granularité};
  \node[uc] (pre) at (10.30, 0.85) {Prétraiter\\le signal};
  \node[uc] (cls) at (10.30,-0.40) {Classer les époques};
  \node[uc] (met) at (10.30,-1.75) {Calculer les\\métriques AASM};

  % ── Cas étendu ──────────────────────────────────────────────────────────
  \node[uc] (osa) at (10.30,-3.65) {Prédire la\\sévérité};

  % ── Acteur ──────────────────────────────────────────────────────────────
  \pic at (1.05,0.55) {acteur};
  \node[nomact,anchor=south] at (1.05,0.97) {Médecin};

  \draw[lien] (1.40,0.35) -- (ana.west);
  \draw[lien] (1.40,0.35) -- (exp.west);
  \draw[lien] (1.40,0.35) -- (pdf.west);

  % ── Relations d'inclusion (libellés étagés le long du faisceau) ──────────
  \draw[incl] (ana) -- node[et,text=hypnoblue!85,pos=0.30] {$\ll$include$\gg$} (chn);
  \draw[incl] (ana) -- node[et,text=hypnoblue!85,pos=0.44] {$\ll$include$\gg$} (pre);
  \draw[incl] (ana) -- node[et,text=hypnoblue!85,pos=0.58] {$\ll$include$\gg$} (cls);
  \draw[incl] (ana) -- node[et,text=hypnoblue!85,pos=0.72] {$\ll$include$\gg$} (met);

  % ── Relations d'extension ───────────────────────────────────────────────
  \draw[ext] (osa) -- node[et,text=hypnored!85,pos=0.5] {$\ll$extend$\gg$} (met);
  \draw[ext] (pdf) -- node[et,text=hypnored!85,pos=0.5] {$\ll$extend$\gg$} (exp);

\end{tikzpicture}
\vspace{0.2cm}
\caption[Raffinement du cas « Analyser un enregistrement »]%
        {Décomposition du cas d'utilisation central en sous-cas inclus et étendus.}
\label{fig:uc-analyse}
\end{figure}

\subsection{Descriptions textuelles}

Un diagramme montre la structure des interactions mais pas leur déroulement. Les
tableaux~\ref{tab:desc-analyse} et~\ref{tab:desc-osa} décrivent textuellement les deux
cas les plus critiques. Les autres suivent le même format et figurent en annexe.

\begin{table}[H]
\centering
\caption{Description textuelle du cas « Analyser un enregistrement »}
\label{tab:desc-analyse}
\small
\begin{tabular}{@{}L{3.0cm}L{11.4cm}@{}}
\toprule
\textbf{Identifiant} & UC-03 \\
\textbf{Acteur principal} & Médecin du sommeil \\
\textbf{Acteurs secondaires} & Moteur d'inférence, service de stockage \\
\textbf{Préconditions} & Le médecin est authentifié ; sa licence est valide ; il
dispose d'un fichier d'enregistrement au format EDF \\
\textbf{Postconditions} & Un examen est créé dans le dossier du patient, la séquence
de stades et les métriques sont enregistrées, le fichier brut est archivé \\
\midrule
\textbf{Scénario nominal} &
\begin{minipage}[t]{11.2cm}
\begin{enumerate}[leftmargin=1.1em,itemsep=1pt,topsep=2pt]
  \item Le médecin sélectionne le patient concerné.
  \item Il choisit le montage disponible, à deux ou cinq dérivations.
  \item Il choisit la granularité de classification, à trois ou cinq stades.
  \item Il dépose le fichier d'enregistrement.
  \item Le système décode le fichier, résout les dérivations et prétraite le signal.
  \item Le moteur d'inférence classe chaque époque et le système calcule les métriques.
  \item Le système affiche l'hypnogramme, les métriques et les alertes cliniques.
  \item Le système crée l'examen et lance l'archivage du fichier en arrière-plan.
\end{enumerate}
\end{minipage} \\
\midrule
\textbf{Scénarios alternatifs} &
\begin{minipage}[t]{11.2cm}
\begin{itemize}[leftmargin=1.1em,itemsep=1pt,topsep=2pt]
  \item \textbf{4a} --- Le fichier n'est pas au format attendu : le système refuse le
        dépôt et le signale.
  \item \textbf{5a} --- Une dérivation est absente : le système applique son repli
        documenté et poursuit.
  \item \textbf{5b} --- L'enregistrement est plus court qu'une époque : le traitement
        s'interrompt avec un message explicite.
  \item \textbf{8a} --- L'archivage échoue : les résultats restent consultables et le
        transfert est signalé en erreur.
\end{itemize}
\end{minipage} \\
\bottomrule
\end{tabular}
\end{table}

\begin{table}[H]
\centering
\caption{Description textuelle du cas « Prédire la sévérité du SAOS »}
\label{tab:desc-osa}
\small
\begin{tabular}{@{}L{3.0cm}L{11.4cm}@{}}
\toprule
\textbf{Identifiant} & UC-04 \\
\textbf{Acteur principal} & Médecin du sommeil \\
\textbf{Acteur secondaire} & Moteur d'inférence \\
\textbf{Préconditions} & Une séquence de stades est disponible, issue d'une analyse ou
d'une saisie manuelle \\
\textbf{Postconditions} & Une classe de sévérité, ses probabilités et la décomposition
des facteurs sont associées à l'examen \\
\midrule
\textbf{Scénario nominal} &
\begin{minipage}[t]{11.2cm}
\begin{enumerate}[leftmargin=1.1em,itemsep=1pt,topsep=2pt]
  \item Le système extrait les marqueurs d'architecture depuis l'hypnogramme.
  \item Le médecin complète les données démographiques et oxymétriques.
  \item Il déclenche l'évaluation du risque.
  \item Le système construit le vecteur de variables et comble les valeurs absentes.
  \item Le moteur produit la classe de sévérité et ses probabilités.
  \item Le système calcule les contributions de chaque variable.
  \item Le système affiche la sévérité, les facteurs aggravants et protecteurs, et
        l'interprétation clinique.
\end{enumerate}
\end{minipage} \\
\midrule
\textbf{Scénarios alternatifs} &
\begin{minipage}[t]{11.2cm}
\begin{itemize}[leftmargin=1.1em,itemsep=1pt,topsep=2pt]
  \item \textbf{2a} --- Le médecin importe un rapport externe : le système en extrait
        les valeurs et préremplit le formulaire.
  \item \textbf{4a} --- Des variables cliniques manquent : le système applique la
        valeur médiane de la population d'apprentissage et le signale.
  \item \textbf{5a} --- Les modèles ne sont pas chargés : l'opération est refusée avec
        un message explicite.
\end{itemize}
\end{minipage} \\
\bottomrule
\end{tabular}
\end{table}

% =============================================================================
\section{Conclusion}
% =============================================================================

Ce chapitre a traduit les intentions des deux chapitres précédents en une commande
vérifiable. Quatre acteurs ont été identifiés, dont deux secondaires qui matérialisent
les frontières techniques du système. Les besoins fonctionnels ont été regroupés en
neuf épics et déclinés en soixante récits utilisateur, priorisés selon MoSCoW et
estimés à 358 points de complexité. Les besoins non fonctionnels ont été énoncés avec,
pour chacun, son moyen de vérification.

L'ordonnancement de ce backlog a produit six sprints de quatre semaines, dont la
vélocité observée — de 29 points au premier à 83 au dernier — reflète la montée en
compétence sur la chaîne d'apprentissage profond plutôt qu'une planification
uniforme. La modélisation des cas d'utilisation a enfin précisé les interactions, en
distinguant les sous-cas systématiquement exécutés des comportements optionnels.

Le chapitre suivant passe de la spécification à la conception : architecture logicielle,
modèle de données, architectures des réseaux de neurones et enchaînements dynamiques.
